% Generated draft: requires mathematical and source review.
\documentclass[11pt]{article}
\usepackage[margin=25mm]{geometry}
\usepackage{fontspec}
\setmainfont{Latin Modern Roman}
\usepackage{amsmath,amssymb,mathtools,unicode-math}
\setmathfont{Latin Modern Math}
\usepackage{enumitem,needspace,xurl,hyperref}
\setlength{\parindent}{0pt}
\setlength{\parskip}{5pt}
\setlength{\emergencystretch}{3em}
\newcommand{\N}{\mathbb{N}}
\newcommand{\Z}{\mathbb{Z}}
\newcommand{\Q}{\mathbb{Q}}
\newcommand{\R}{\mathbb{R}}
\newcommand{\C}{\mathbb{C}}
\newcommand{\F}{\mathbb{F}}
\newcommand{\A}{\mathbb{A}}
\newcommand{\PP}{\mathbb{P}}
\newcommand{\E}{\mathbb{E}}
\newcommand{\eps}{\varepsilon}
\newcommand{\dd}{\,\mathrm{d}}
\begin{document}
\section*{UnsolvedMath ID: 701132}
\textbf{Source problem code:} catalogue entry
\subsection{Definitions and mathematical statement}
Let \(K\) be an algebraically closed field of arbitrary characteristic, let \(G\) be a reductive algebraic group over \(K\), and fix a Borel subgroup \(B\leq G\). For an affine \(G\)-variety \(X\), write \(A=K[X]\) for its coordinate algebra and \(A^G=\{f\in A:g\cdot f=f\text{ for every }g\in G\}\) for its invariant ring. A commutative Noetherian ring \(R\) is Cohen--Macaulay if \(\operatorname{depth}R_{\mathfrak p}=\dim R_{\mathfrak p}\) for every prime ideal \(\mathfrak p\).

\textbf{Part (a).} Suppose that the rational \(G\)-module \(A\) has an exhaustive ascending filtration by \(G\)-submodules whose successive quotients have the form \(\operatorname{Ind}_B^G V\), where each \(V\) is a one-dimensional rational \(B\)-module and may vary with the quotient. Here \(\operatorname{Ind}_B^G\) denotes rational induction from \(B\) to \(G\). Suppose also that \(A\) is Cohen--Macaulay. Must \(A^G\) be Cohen--Macaulay?

The source additionally describes the Cohen--Macaulay hypothesis by saying that \(A\) is free over the subalgebra generated by any homogeneous system of parameters. However, it specifies neither a grading on \(A\) nor hypotheses under which this description is equivalent to the usual ring-theoretic definition above. This qualification remains unresolved; the displayed question uses the usual definition of Cohen--Macaulayness. The supplied comment dated 2005 reports an affirmative result when \(X\) is a rational \(G\)-module, citing M. Hashimoto.

\textbf{Part (b), archived separately as solved.} The source asks whether \(K[M(n)^m]^{GL(n)}\) is Cohen--Macaulay in all characteristics, where \(M(n)^m\) is the direct sum of \(m\) copies of the space of \(n\times n\) matrices. It does not explicitly specify the action or the ranges of \(n\) and \(m\). Under the customary interpretation of simultaneous conjugation and positive integer parameters, this asks: for every algebraically closed field \(K\) of arbitrary characteristic and every \(n,m\in\mathbb Z_{\geq1}\), is
\[
K[\operatorname{Mat}_n(K)^{\oplus m}]^{\operatorname{GL}_n(K)}
\]
Cohen--Macaulay, where
\[
g\cdot(T_1,\ldots,T_m)=(gT_1g^{-1},\ldots,gT_mg^{-1})?
\]
This explicit formulation of part (b) is conditional on those interpretations; they require source confirmation. The supplied archive records an affirmative answer attributed to M. Hashimoto, rather than supplying a proof here.
\subsection{Short English statement}
Part (a) asks whether taking invariants preserves Cohen--Macaulayness for an affine reductive-group variety whose coordinate algebra has a filtration by modules induced from one-dimensional Borel-subgroup modules. The source reports the rational-module case as proved. Part (b), supplied as archived solved, concerns invariants of tuples of matrices; its usual simultaneous-conjugation interpretation is stated conditionally because the action and parameter ranges are omitted. The grading needed for the source's homogeneous-parameter description also remains unspecified.
\subsection{Sources}
\begin{itemize}
\item Supplied local UnsolvedMath registry record, permanent ID \(701132\), exact source problem code \texttt{catalogue entry}, title \emph{Kourovka Notebook Problem 15.44}. Proposed by A. N. Zubkov, with proposed year 2002.
\item \emph{The Kourovka Notebook (Unsolved Problems in Group Theory)}, 15th issue (2002), Problem 15.44; the supplied background identifies an update in the 21st issue (2026), \texttt{arXiv:1401.0300v48}. The supplied collection description names E. I. Khukhro and V. D. Mazurov as editors.
\item S. Donkin, \emph{Rational representations of algebraic groups. Tensor products and filtration}, Lecture Notes in Mathematics, vol. 1140, Springer, Berlin, 1985. Supplied reference for the representation-theoretic filtration in part (a).
\item M. Hashimoto, \emph{Math. Z.}, vol. 236 (2001), pp. 605--623. The supplied 2005 comment cites this reference for the rational \(G\)-module case of part (a); the supplied Archive of Solved Problems entry records an affirmative answer for the archived part (b) with the same reference. No article title is supplied.
\item The local record labels the problem \texttt{partially\_solved}. Its status review, checked on \texttt{2026-10-16}, states that solved parts occur in the archive and remaining parts in the main section. This is supplied dated curation, not an independently verified mathematical or literature assessment.
\end{itemize}
\end{document}
